\section{Conclusion}
\label{sec:conclusion}

Starting from stock MACE-OFF23 medium in FP64, nine measured steps brought a complete energy+force evaluation on one H100 to 0.81\,ms at 21 atoms, 3.0\,ms at 648 atoms and 23.9\,ms at 5{,}184 atoms in a periodic box (12--20$\times$ at the complete-step boundary), with zero host synchronisations per MD step, forces within $2\times10^{-14}$\,\eVA{} of stock, and NVE drift indistinguishable from stock. The largest single factors were CUDA-graph capture for small systems (12$\times$ at aspirin), and the node-centric convolution plus GPU cell list for large ones (2.2$\times$ and 2.8$\times$). At $\ge$2{,}000 atoms the step is now $\sim$99\% GPU-busy inside the model, so further progress has to come from the model's kernels (convolution backward, radial-MLP GEMMs and elementwise glue, together 75\% of the time) or from relaxing FP64.

The most transferable outputs are probably not the speedups but the traps: an FP64 library mode that is FP32-accurate while passing its own tests, CUDA-graph captures that succeed without capturing, graph reuse keyed only on shape, padding placement that serialises a destination-owned kernel, and floating-point evaluation orders that must be matched for bitwise neighbour lists. Before a state-of-the-art claim is made, the result must be compared against MACE's own cuEquivariance + \code{torch.compile} path at the same boundary.
